\documentclass[a4paper,12pt,]{article}

\usepackage{amsmath,amssymb}
\usepackage{listings}
\usepackage{xcolor}

\lstset{
	language=Python,
	basicstyle=\ttfamily\small,
	backgroundcolor=\color{gray!10},
	frame=single,
	numbers=left,
	numberstyle=\tiny,
	breaklines=true,
	showstringspaces=false
}

\title{Übungsblatt 1.2}
\author{Alexander Engelhardt}
\date{\today}

\begin{document}
	\maketitle
	
	\section*{Aufgabe 1.6}
	\subsection*{Teilaufgabe 1.6.2}
	\subsubsection*{a)}
	\begin{enumerate}
		\item Die \textbf{gemessenen Höchstgeschwindigkeiten} der Stichprobe sind durchschnittlich 65.74 km/h vom arithmetischen Mittel (289 km/h) entfernt.
		\item Unter 373 km/h (dem 90\%-Quantil) liegen \textbf{90\% aller gemessenen Werte.}
	\end{enumerate}
	\subsubsection*{b)}
	
	\subsubsection*{c)}
	\section*{Aufgabe 1.7}
	\subsection*{Teilaufgabe 1.7.1}
	\paragraph{Der \textbf{Korrelationskoeffizent -0.966} ist zu der Punktwolke \textit{(B)} gehörig.}
	\begin{enumerate}
		\renewcommand{\labelenumi}{(\Alph{enumi})}
		\item Der Zusammenhang von diesen Punkten und des Achsenkreuzes ist stets positiv linear, und somit, ipso facto, kann er nicht in solch negativer Latenz sein.
		\item 
		\begin{lstlisting}
Bc = np.corrcoef(X, By)
print("Korrelationskoeffizient von:\nB: ", Bc) 
# [[ 1.         -0.96609178]
# [-0.96609178  1.        ]]
		\end{lstlisting}
		\item Ist am Anfang zu konstant, um solch einen negativen Korrelationskoeffizienten zu besitzen, selbst beim (von Punkt 4 aus) rapiden Abfall von Punkt 5 und Punkt 6.
	\end{enumerate}
	\subsection*{Teilaufgabe 1.7.2}
	\begin{enumerate}
		\item R = A: 0.33
		\item R = B: 0.69
		\item R = C: 0.98
		\item R = D: 1.00
		\item R = E: 0.00
		\item R = F: -0.64
		\item R = G: -0.92
		\item R = H: -1.00
	\end{enumerate}
	\subsection*{Teilaufgabe 1.7.3}
	\begin{enumerate}
		\item $\rho = K: 1.0$
		\item $\rho = L: 0.3$
		\item $\rho = M: -0.7$
		\item $\rho = N: 0.9$
		\item $\rho = O: -0.3$
		\item $\rho = P: -0.9$
		\item $\rho = Q: 0.7$
		\item $\rho = R: 0.0$
		\item $\rho = S: -1$
	\end{enumerate}
	\section*{Aufgabe 1.10}
	\subsection*{Teilaufgabe 1.10.1}
	\paragraph{\textbf{Qualitativ, Quantitativ, Diskret, Stetig:} Anzahl der Räder des Fahrzeugs: Dieses Merkmal ist \textbf{diskret}, da ein Fahrzeug nur ganze Reifen haben kann.}
	\subsection*{Teilaufgabe 1.10.2}
	
\end{document}